\documentclass[a4paper,11pt]{article}

\usepackage{times}
\usepackage[T1]{fontenc}
%\usepackage[xetex]{graphicx} % when you compile with xelatex
\usepackage[pdftex]{graphicx} % when you compile with pdflatex
\usepackage{epstopdf}
\usepackage{float}
\usepackage{enumerate}
\usepackage{boxedminipage}
\usepackage{url}
\usepackage{hyperref}
%\usepackage[dutch]{babel} % when you write in dutch

\usepackage{dcolumn}
\newcolumntype{d}[1]{D{.}{.}{#1}}

%%%%%%%%%%%% layout %%%%%%%%%%%%%%%%%%
\usepackage[margin=3cm]{geometry}
% you might enhance readability by increasing the distances between lines:
%\renewcommand{\baselinestretch}{1.1}

%%%%%%%%%%% mathematics %%%%%%%%%%%%%%%%%%
\usepackage{amsmath,amsfonts,amssymb}

%%%%%%%%%%%%% theorems %%%%%%%%%%%%%%%%
% adapt when you write Dutch

\usepackage{amsthm}

\theoremstyle{plain}
\newtheorem{theorem}{Theorem}
\numberwithin{theorem}{subsection}
\newtheorem{corollary}{Corollary}
\numberwithin{corollary}{subsection}
\newtheorem{proposition}{Proposition}
\numberwithin{proposition}{subsection}
\newtheorem{lemma}{Lemma}
\numberwithin{lemma}{subsection}
\newtheorem{assumption}{Assumption}
\numberwithin{assumption}{subsection}

\theoremstyle{definition}
\newtheorem{definition}{Definition}
\numberwithin{definition}{subsection}
\newtheorem{example}{Example}
\numberwithin{example}{subsection}
\newtheorem{remark}{Remark}
\numberwithin{remark}{subsection}
\newtheorem{notation}{Notation}
\numberwithin{notation}{subsection}

%%%%%%%%%%%%algorithms %%%%%%%%%%%%%%%%%%%
\usepackage{algorithm}
\usepackage{algorithmic}

%%%%%%%%% bibliography style %%%%%%%%%%%%%
\usepackage{natbib} % in case of classic bibtex
%\usepackage[natbib,style=authoryear,maxcitenames=2,
%           maxbibnames=99,hyperref,backend=biber]{biblatex}
%\addbibresource{bib/EORthesis.bib}

%%%%%%%%%%%%%%%% start document %%%%%%%%%%%
\begin{document}

%%%%%%%%%%%%%%% front page %%%%%%%%%%%%%%%%

\thispagestyle{empty}

\includegraphics[height=2cm]{figures/LogoSBE.png}

\vspace*{3cm}

\noindent
\rule{\textwidth}{0.8pt}
\begin{center}
{\huge\bf
\noindent
Deriving fractional moments using the Moment Generating Function
}
\end{center}

\vspace*{-8pt}
\noindent
\rule{\textwidth}{0.8pt}

\vspace*{2cm}

\begin{center}
{\LARGE\bf
Jelle Reisinger
}

{\Large
\vspace*{0.5cm}
(2780350)


\vspace*{2cm}

30 June 2025
}
\end{center}

\vspace*{2cm}

{\Large
\noindent
Bachelor Thesis Econometrics and Data Science
}

\vspace*{1cm}

{\Large
\noindent
Thesis commission:\\[0.3cm]
PhD candidate Gabriele Mingoli\\[0.3cm]
Dr. yyyy (co-reader)
}


\newpage

%%%%%%%%%%%%%%%%%%%%%%% main contents %%%%%%%%%%%%%%%%%
\setcounter{page}{1}

\section*{Abstract}
The abstract should summarize the contents of the thesis.
It should be clear, descriptive, self-explanatory and not longer
than a third of a page. Please avoid using mathematical
formulas as much as possible.
Keywords might be given.

\bigskip\noindent
\textbf{Keywords:} Fractional moments, Moment Generating Function, fractional calculus.

\section{Introduction}\label{s:intro}
Statistical moments such as the mean and the variance are essential tools to characterize data and its distribution. Moments of even higher order are useful regarding the shape of the distribution. Less known moments, however, are the fractional moments. While these moments may not  be that significant in comparison with the first for integer moments of a distribution, they can be very useful in certain applications. While these moments can be obtained by ordinary integration, the field of fractional calculus, which finds many application in physics, provides a more general framework to compute these moments. In this paper, we will combine the aforementioned technique with the Moment Generating Function to obtain new expressions for the fractional moments of a distribution.
\section{Literature Review}
\subsection{Application of fractional moments}
Fractional moments play a significant role in a variety fields, including finance, economics, and statistics An example of the latter is its application in approximating integer moments as described by \cite{inverardi2024}. This is especially useful when (for \(\alpha \in \mathbb{R}\)) the \(\lceil\alpha \rceil\)-th central integer moment might not exist, while its fractional moment does. Finding an existing fractional moment close to a non-existing integer moment, could still provide one with information about this integer moment. For example, the student-t distribution with \(\nu = 2\) degrees of freedom only has central and raw moments of order \(k\), where \( 0 < k < \nu\). This implies that its second central moment \((k = 2)\) does not exist. One could however consider the \(k-th\) central moment where \(k = 1.95\) and interpret its value as the variance of the distribution. Fractional moments are also used in financial modelling, particularly in the context of Generalized Autoregressive Conditional Heteroskedasticity (GARCH) models. These models are commonly applied to time series data such as financial returns and capture the dynamic volatility that changes over time. The GARCH model achieves this by modeling the volatility based on the returns and variances of previous time periods. \cite{hansen2024} have obtained a method of finding fractional absolute moments of the cumulative return, which would have been impossible when using any other method. \cite{gyzl2013} have also introduced the usage of Fractional moments in risk-models, specifically insurance models. In such models, often the probability density function of total ruin, the event where an insurance company's capital becomes negative, is unknown. The so-called "Method of maximum entropy" has been developed to find these densities. This method takes fractional moments as its input, as they have been proven to be able to characterize its distribution \cite{lin1992}. This method has proven to be a useful alternative to existing methods, such as inverse Laplace transformations. This is the case as this new method takes less values than the latter as input, making it computationally more efficient. Beyond finance and risk modelling, fractional moments also have important applications in engineering. Examples include optimizing signal processing and control systems as well as studying the response characteristics of random vibration systems. \cite{wang2025} has shown that when using the concept of fractional moments for the latter, accuracy and stability is higher compared to traditional methods, such as Taylor expansions. Furthermore, in terms of simplicity and efficiency, the method of fractional moments is advantageous, as its computation steps are straightforward and avoid convergence issues, significantly reducing the resources required for computation. Working with fractional moments has allowed \cite{wang2025} to obtain both analytical, as well as numerical solutions to problems within their research field, which again proves its viability. Another application within the field of engineering, can be found in the identification of distributions of non-linear systems. \cite{dimatteo2014} have shown that complex fractional moments allow one to solve equations such as the Kolmogorov or Fokker-Planck equation, a characterization of continuous-time Markov processes. After performing a Mellin transformation on this system of equations, the resulting system is a linear system in terms of complex fractional moments. The latter can now be solved rather easily and taking the Inverse Mellin transformation on these solutions immediately provides one with the solutions of the non-linear system. Advantages of using complex fractional moments instead of integer moments is that, when applying the Mellin transformation, the relevant PDF is restored on its entire support. This is not the case for the integer moments. This method of using complex fractional moments has been proven to have a rather high accuracy and is applicable to any stable kind of non-linear system of equations \cite{dimatteo2014}.

\subsection{Obtaining fractional moments}
The traditional method of computing fractional moments is rather straightforward. Similarly to integer moments, one simply computes the summation (in the discrete case) or integral (in the continuous case) of \(x^\alpha \cdot f_x(x)\), where \(f_x(x)\) denotes the probability density function of the random variable \(X\). In the context of fractional moments, \(\alpha \in \mathbb{R}\) instead of \(\mathbb{Z}\) (assuming that negative moments exist). \cite{hansen2024} introduce an alternative approach to computing fractional moments using the complex moment generating function (CMGF), which they apply in the context of the aforementioned GARCH models. One of their key expressions is given by:

\[\mathbb{E}\left| X - \xi \right|^r = \frac{\Gamma(r+1)}{2\pi} \int_{-\infty}^{\infty} \frac{e^{-\xi z} M_X(z) + e^{\xi z} M_X(-z)}{z^{r+1}} dt\] \(\text{ where } z = s + it, s \in \mathbb{N_+}, \xi \in \mathbb{R} \) and \(r\) of course the order of the moment.

This formulation extends upon the traditional moment generating function (MGF) but avoids the process of taking derivatives, making it computationally efficient. The inclusion of the Gamma function is logical, as it extends the factorial function to real values, aligning well with the computation of fractional moments. Since this method relies on integration, rather than differentiation, it avoids numerical issues that might arise when computing derivatives, such as obtaining rather great approximation errors.
While the CMGF method provides an efficient  and elegant alternative, this thesis explores a different approach: computing fractional moments directly by applying fractional derivatives to the MGF. The MGF is widely used for computing integer moments by differentiation around zero. To extend this approach to fractional orders requires us to take fractional derivatives. Thus, we need to define such fractional derivative operators. These fractional derivatives have a long history and often make use of the aforementioned Gamma function in combination with some integral. This means that, for continuous random variables, where we integrate the MGF, we will have to do double integration. A consequence might be that obtaining analytical expressions of these moments may not always be possible. A lot of alternative expressions of these fractional derivatives have been created, mostly based on different interpretations of the latter in the field of physics. This implies that different expressions of fractional derivatives in combination with the MGF might obtain different  moments expressions for the same distribution and same fractional order of the moment. Thus, it is essential to compare each of these definitions with the traditional way of computing fractional moments, to derive their accuracy and conclude which approach is most suitable for fractional moment computation. Similar to the expression of the moments of a random variable, their 'Biases' may also be hard to derive analytically, depending on its distribution.



\subsection{Fractional derivatives}
Although not the main topic of interest of this thesis, it is useful to have some knowledge about the history and applications of fractional derivatives. The study of fractional derivatives has been relevant as early as the year 1695 when the concept of such a derivative was implicitly discussed by Leibniz and Bernoulli \cite{katugampola2014}. Since then, numerous definitions of fractional derivatives have been developed. The best-known definition of the fractional derivative is the Riemann-Liouville derivative, its upper derivative of a function \(f(x)\) of order \(\alpha\) is denoted as 
\[\frac{d^n}{dx^n}\frac{1}{\Gamma(n -\alpha)} \int_{a}^{x} (x-t)^{n - \alpha -1} f(t) dt, \text{ where }n = \lceil\alpha \rceil \] \cite{kilbas2006}.
Michele \cite{caputo1967} defined a variation on this derivative, where instead of \(\frac{d^n}{dx^n}\) in front of the integral, we have \(\frac{d^n}{dt^n}\) inside the integral. Due to this adjustment, it is possible to have initial value conditions expressed as the traditional derivatives of integer-order, which made these fractional differential equation problems more intuitive. Other fractional derivatives, such as the \cite{hadamard1892} and \cite{riesz1949} derivative, have been defined to take advantage of particular beneficial properties. For example, each of the latter derivatives can be written as a Fourier transformation. As a consequence, the analytical expressions, can often be simplified. A rather unique derivative is the Grünwald-Letnikov derivative which, in contrast to all the aforementioned derivatives, is not based on integral. Instead, it generalizes the difference quotient, \(\frac{f(x+h) - f(x)}{h}\),  to fractional orders using binomial coefficients \cite{atici2021}. This variety of definitions emphasises how dependent fractional derivatives are  on different physical interpretations and practical applications.
Beyond their theoretical significance, fractional derivatives have been of significant importance in various scientific fields since the 19th century. Examples include fractional Fourier transformations, a generalization of the regular Fourier transformations \cite{missbauer2012}, fractional diffusion equation models, describing the motion of particles in liquids as a consequence of thermal molecular motions \cite{einstein1905} and the fractional Schrödinger equation, a generalization of the Schrödinger equation, often used in quantum mechanics \cite{laskin2002}. Their applications are less common in the fields of finance or economics, as fractional derivatives are mainly used to describe natural phenomena \cite{boulaaras2023}. Yet they still offer some great potential. (Symmetric) Levy flights make use of fractional derivatives in order to solve partial differential equations which describe random walk processes in time series \cite{scalas2000}. The development of fractional derivatives also led to the notion of fractional Brownian motions, a generalization of the Brownian motion \cite{mandelbrot1968}. The latter is a continuous-time stochastic process which, similar to Levy flights, may be used to model random walk processes.
\section{Fractional Calculus}\label{s:calculus}
\subsection{Overview and applications of fractional derivatives}\label{ss:calculus_introduction}
Although not the main topic of interest of this thesis, it is useful to have some knowledge about the history and applications of fractional derivatives. The study of fractional derivatives has been relevant as early as the year 1695 when the concept of such a derivative was implicitly discussed by Leibniz and Bernoulli \cite{katugampola2014}. Since then, numerous definitions of fractional derivatives have been developed. The best-known definition of the fractional derivative is the Riemann-Liouville derivative, its upper derivative of a function \(f(x)\) of order \(\alpha\) is denoted as 
\[\frac{d^n}{dx^n}\frac{1}{\Gamma(n -\alpha)} \int_{a}^{x} (x-t)^{n - \alpha -1} f(t) dt, \text{ where }n = \lceil\alpha \rceil \] \cite{kilbas2006}.
Michele \cite{caputo1967} defined a variation on this derivative, where instead of \(\frac{d^n}{dx^n}\) in front of the integral, we have \(\frac{d^n}{dt^n}\) inside the integral. Due to this adjustment, it is possible to have initial value conditions expressed as the traditional derivatives of integer-order, which made these fractional differential equation problems more intuitive. Other fractional derivatives, such as the \cite{hadamard1892} and \cite{riesz1949} derivative, have been defined to take advantage of particular beneficial properties. For example, each of the latter derivatives can be written as a Fourier transformation. As a consequence, the analytical expressions, can often be simplified. A rather unique derivative is the Grünwald-Letnikov derivative which, in contrast to all the aforementioned derivatives, is not based on integral. Instead, it generalizes the difference quotient, \(\frac{f(x+h) - f(x)}{h}\),  to fractional orders using binomial coefficients \cite{atici2021}. This variety of definitions emphasises how dependent fractional derivatives are  on different physical interpretations and practical applications.
Beyond their theoretical significance, fractional derivatives have been of significant importance in various scientific fields since the 19th century. Examples include fractional Fourier transformations, a generalization of the regular Fourier transformations \cite{missbauer2012}, fractional diffusion equation models, describing the motion of particles in liquids as a consequence of thermal molecular motions \cite{einstein1905} and the fractional Schrödinger equation, a generalization of the Schrödinger equation, often used in quantum mechanics \cite{laskin2002}. Their applications are less common in the fields of finance or economics, as fractional derivatives are mainly used to describe natural phenomena \cite{boulaaras2023}. Yet they still offer some great potential. (Symmetric) Levy flights make use of fractional derivatives in order to solve partial differential equations which describe random walk processes in time series \cite{scalas2000}. The development of fractional derivatives also led to the notion of fractional Brownian motions, a generalization of the Brownian motion \cite{mandelbrot1968}. The latter is a continuous-time stochastic process which, similar to Levy flights, may be used to model random walk processes.
\subsection{Formal definitions of fractional derivatives}
In order to obtain the expressionS mentioned in \autoref{s:methodology}, some advanced tools are required. We can find these in the field of fractional calculus.
We define the following:
\begin{definition}
    Let \(D\) be the differential operator, such that \(D f(x) = \frac{d}{dx} f(x)\). Then the fractional derivative of order \(\alpha\) is defined as \[D^{\alpha} f(x) = \frac{d^{\alpha}}{dx^{\alpha}} f(x)\].
\end{definition}
In this definition, \(\alpha\) can be any real number. When taking regular derivatives, \(\alpha \in \mathbb{N}\). In most of our cases, we are interested in the instance where  \(\alpha \in \mathbb{R}_+\).

It is also possible to study derivatives of negative order, which can be used to obtain moments of negative order of a function, provided that such an order exists. A derivative of negative order is simply an integral of positive order. This is defined as follows:
\begin{definition}
    Let \(I\) be the integral operator, such that \(I f(x) = \int f(x) dx\). Then the fractional integral of order \(\alpha\) is defined as \[(I^{\alpha} f) (x) = \frac{1}{(\alpha-1)!}\int (x-t)^{\alpha-1} f(t) dt\] \cite{cauchy1823}.
\end{definition} 

Combining the previous two definitions, we obtain the following, more general, definition.
\begin{definition}\label{d: differintegral}
    The differintegral operator is defined as
    \begin{equation}
        R^\alpha f(x) = \begin{cases}
            I^{|\alpha|} f(x) & \text{if } \alpha < 0 \\
            D^\alpha f(x) & \text{if } \alpha > 0 \\
            f(x) & \text{if } \alpha = 0
        \end{cases}, \text{ with } \alpha \in \mathbb{R}.
        \end{equation}
\end{definition}

As mentioned in section \ref{ss:calculus_introduction}, quite a number of different definitions have been proposed to compute a fractional derivative. Some of these definitions are rather similar, thus, in this paper, we will focus on some of the more well known fractional derivatives. We will start with the most famous fractional derivative, which laid the foundation of the study of fractional derivatives as early as in 1832.
\begin{definition}
    The left-side Riemann-Liouville fractional derivative of order \(\alpha\) is defined as:
    \begin{equation}
        D^{\alpha}_{a_+} f(x) =  \frac{d^{n}}{dx^{n}} D_{x}^{-(n - \alpha)} f(x) = \frac{d^{n}}{dx^{n}} I_{x}^{n - \alpha} f(x) = \frac{d^n}{dx^n} \frac{1}{\Gamma(n -\alpha)}  \int_{a}^{x} (x-t)^{n - \alpha-1} f(t) dt
    \end{equation} \cite{liouville1832}.

    Here, \(n = \lceil \alpha \rceil\), the ceiling function and \(\Gamma(.)\) is the Gamma function, see section \autoref{s:appendices}.

    Note that we just defined the left-side Riemann-Liouville fractional derivative, suggesting that there also exists a right side derivative. In the case of the latter, we would evaluate the associated integral the other way around. Namely, \(D^{\alpha}_{b_-} f(x) = \frac{d^n}{dx^n} \frac{1}{\Gamma(n -\alpha)}  \int_{x}^{b} (x-t)^{n - \alpha-1} f(t) dt\). We intend on using the left-side derivative, as is supported by \cite{tarasov2023}. The reason is, due to the fact that many functions in probability theory, most importantly the cumulative distribution function, are defined as an integral from some constant to \(x\). 
    
    \begin{remark}\label{r: integer}
        For values \(\alpha \in \mathbb{N}_+, n =  \lceil \alpha \rceil = \alpha\), so \(\Gamma(n - \alpha) = \Gamma(0)\), which is undefined. Thus for \(\alpha \in \mathbb{N}_+\), we define: \(D^\alpha f(x) = \frac{d^\alpha}{dx^\alpha} f(x)\), which is simply the regular expression for derivatives of integer order.
    \end{remark}
   
\end{definition}
A modification of the Riemann-Liouville derivative is the Caputo-Fabrizio derivative, which is defined as follows:
\begin{definition}\label{d: CF}
    The Caputo-Fabrizio fractional derivative of order \(\alpha, \alpha \in [0,1)\) is defined as:
    \begin{equation}
        D^{\alpha} f(x) = \frac{1}{1 - \alpha}  \int_{a}^{x} \exp\left(\frac{-\alpha}{1 - \alpha}(x-t)\right) f'(t) dt
    \end{equation} \cite{caputo2015}.
    With \(a \in [-\infty, x)\).
    The Caputo-Fabrizio is always defined as the integral from some constant to the variable \(x\). This is another reason for choosing to work with the left-side Riemann-Liouville integral. In this way, it will be more straightforward to compare the two integrals. What is more, note that for the Caputo-Fabrizio derivative, the order \(\alpha \in [0, 1)\). This does not mean, however, that one can only compute fractional derivatives of order 1 or lower. There exists a rather convenient property of the differintegral operator which allows one to combine orders of derivatives, which will be discussed in a moment.
    
\end{definition}

Lastly, we will define the Grünwald-Letnikov derivative, which is defined as follows:
\begin{definition}
    The Grünwald-Letnikov fractional derivative of order \(\alpha\) is defined as:
    \begin{equation}
        D^\alpha f(x) = \lim_{h \to 0} \frac{1}{h^\alpha} \sum_{k=0}^\infty (-1)^k \binom{\alpha}{k} f(x - k h)
    \end{equation}
   Where \(\binom{\alpha}{k}\) is the binomial coefficient, with \(0 \leq k \leq \alpha\), see section \autoref{s:appendices}.
\end{definition} \cite{zhmakin2022}.
It is immediately clear, observing the summmation symbol instead of the integral, that this derivative behaves quite differently from the two derivatives defined above. The Grünwald-Letnikov derivative is an extension on derivatives based of the concept of finite differences \cite{flajolet1995}.

We will consider a number of properties which come in useful when working with fractional derivatives.

\begin{proposition}\label{p: calculus}
    The fractional derivatives above adhere to the following properties:
    \begin{enumerate}[(i)]
        \item Linearity: Let \(f(x), g(x)\) be functions and \(a, b, x \in \mathbb{R}\). Then we have that \(D^{\alpha} (a f(x) + b g(x)) = a D^{\alpha} f(x) + b D^{\alpha} g(x)\).
        \item \(D^{\alpha} f(x) = f(x)\), for \(\alpha = 0\) 
        \item for sufficiently smooth functions f, we have that \(D^{\alpha + \beta} f(x) = D^\alpha(D^\beta f(x)) =  D^\beta(D^\alpha f(x))\), with \(\alpha, \beta \in \mathbb{R}\). Note that, for definition \ref{d: CF}, this property only holds for \(\beta \in \mathbb{N}, \alpha \in [0,1)\).
    \end{enumerate}
        
    
\end{proposition}

The proofs of these of the properties stated in this proposition can be found in appendix \ref{s:app_B}. Most of these proofs have been provided by myself, while some other proofs, which are outside the scope of thesis are based on other papers. Note that the third property is especially useful for the Caputo-Fabrizio derivative. This property allows one to take fractional derivatives of order greater than 1 by comparing fractional derivatives and regular integer derivatives.

We will now provide two numerical examples of these fractional derivatives. For simplicity, we will let \(a = 0\):
\begin{example}
    \begin{enumerate}[(i)]
        \item 
    
    We consider the Riemann-Liouville derivative of order \(\frac{3}{2}\) for some constant \(c \in \mathbb{R}\):
    \[D^{\frac{3}{2}}_{RL}(c) = \frac{d^2}{dx^2} \frac{1}{\Gamma(2 - \frac{3}{2})}  \int_{0}^{x} (x-t)^{2 - \frac{3}{2}-1} c dt\]
    \[= \frac{d^2}{dx^2} \frac{c}{\sqrt{\pi}}  \int_{0}^{x} (x-t)^{- \frac{1}{2}} dt = \frac{d^2}{dx^2} \frac{-2c}{\sqrt{\pi}} \sqrt{x - t} \Big|_{0}^{x}\]
    \[= \frac{d^2}{dx^2} \frac{2c \sqrt{x}}{\sqrt{\pi}} = \frac{-c}{2x\sqrt{\pi x}} \neq 0.\] As stated, the fractional derivative of a constant is not equal to zero when using the Riemann-Liouville derivative. This is also the case for the Grünwald-Letnikov derivative, but not for the Caputo-Fabrizio derivative.
    \item We compute the semi-derivative of \(\frac{x}{2}\) using the Caputo-Fabrizio derivative:
    \[D^{\frac{1}{2}}_{CF}\left(\frac{x}{2}\right) = \frac{1}{1 - \frac{1}{2}}  \int_{0}^{x} \exp\left(\frac{-\frac{1}{2}}{1 - \frac{1}{2}}(x-t)\right) \frac{1}{2} dt = \int_{0}^{x} \exp(t - x) dt\]
    \[ =  \exp(t - x) \Big|_{0}^{x} = 1 - \exp(- x).\]
    For \(x \geq 0\), this expression is equal to the CDF of the exponential distribution with \(\lambda = 1\). A remarkble result.
    
    \end{enumerate}
\end{example}
The observant reader might notice that no explicit examples of the Grünwald-Letnikov derivative have been provided. The latter is due to the fact that is rather difficult to obtain analytical expressions for this derivative. Thus, later on in this thesis, when computing fractional moments and their assiocated biases, the main focus for the Grünwald-Letnikov derivative will be on its numerical computations.
\begin{remark}
    
As shortly mentioned in the introduction, it is possible to generalize the order of derivatives even further, extending \(\alpha\) to be in \(\mathbb{C}\) instead of \(\alpha \in \mathbb{R}\). This means that, when combining such derivatives with the moment generating function, we will obtain complex moments. Since statisical moments of complex order do almost not find any usage in applications, as they lack interpretability, they are not the main focus of this research. For the interested reader, \cite{love1971} has done some impressive research on the fundamentals of derivatives of complex order. The obtained expressions are somewhat similar to those of the fractional derivatives which have been discussed above.
\end{remark}
\subsection{The Moment Generating Function of fractional order}\label{s:MGF}
Having defined the techniques required to compute fractional derivatives, we now try to combine them with the MGF. Before proceeding, we will briefly review the formal definition of a moment in statistics.
\subsubsection{Moments}
Recall the definition of statistical moments as in definition \ref{eq:moments} of section \ref{ss:methodology_introduction}. If \(c = \mu_x\), where \(\mu_x\) denotes the expected value of \(X\), then our higher moments are called central moments. In the context of this research, we will focus on the case \(c = 0\), corresponding to raw moments of a random variable \(X\). This choice has been made as the MGF, which we will soon define, only computes raw moments of higher order. A moment of order \(\alpha\) is said to exist, if \(\mathbb{E}[X^\alpha] < \infty\).

\subsubsection{The Moment Generating Function}
We now formally introduce the moment Generating Function, one of the most significant subjects of this thesis.
\begin{definition}
    The moment generating function of a variable \(X\), is defined as
    \[M_X(t) = \mathbb{E}[e^{tX}]\] provided that \(\mathbb{E}[e^{tX}] < \infty\), for all  \( t \in (- h, h)\), which contains 0, for some \(h > 0\) 
\end{definition}

\begin{remark}
    Deriving the expression \(M_X(t) = \mathbb{E}[e^{tX}]\) is typically straightforward. Generally, it simply requires a handful of steps of analytic evaluation. This procedure is not that interesting nor relevant to this research. Thus, when making use of  expressions of the MGF, we will simply refer to the distribution table in Appendix \ref{s:app_common_distributions}.
\end{remark}

 
We will state the theorem which makes the MGF so useful. This theorem allows us to compute moments of higher order by taking derivatives of the given order instead of integrals.
\begin{theorem}\label{t:mgf}
    If \(M_X(t)\) exists on some interval \((-h, h)\), as defined before, we have that:
    \[ \mathbb{E}[X^n] = M_X^{(n)}(0), \text{ for } n \in \mathbb{N}\] 
\end{theorem} 




We introduce the following well-known properties for the MGF \(M_X(t)\):
\begin{proposition}\label{p: moments}
    For \(X, Y\) random variables, we have that:
    \begin{enumerate}[(i)]
        \item \(M_X^{(0)}(t) = \mathbb{E}[e^{0X}] = 1\). This property can be used to confirm that a given function is a valid probability density function (i.e., integrates to one).
        \item Location scale-transform. Assuming \(M_X(t)\) exists, for constants \(\mu, \sigma \in \mathbb{R}\), we have that: 
        \[M_{\mu + \sigma X}(t) = e^{\mu t} \cdot M_X(\sigma t)\]
        \item If \(X \perp Y\), then \(M_{X+Y}(t) = M_X(t)\cdot M_Y(t)\).
    \end{enumerate}
\end{proposition}


The proofs of the latter can be found in Appendix \ref{pf:MGF}.

There are several additional topics closely related to the MGF, including Fourier transforms, Laplace transforms, Wick rotations, and characteristic functions. While these subjects are relevant to the theoretical foundation of the MGF , they fall outside the scope of this thesis and will therefore not be discussed. Readers interested in exploring these concepts further may find \citet{kolmogorov1999} to be a valuable resource.

\subsubsection{Computing moments of negative order using the Moment Generating Function}
In specific cases, as were mentioned in section \ref{s:intro}, moments of negative order can be of interest to characterize the data. Such an example was the moment of order \(-\frac{1}{2}\), which in some contexts may be used to obtain the Sharpe Ratio. Thus, it is useful to understand how to compute moments of such orders. We will consider the continuous case:

\[\mathbb{E}[X^{-n}] = \int_{-\infty}^{\infty} x^{-n} f_X(x) dx = \int_{-\infty}^{\infty} \left(\frac{1}{x}\right)^n f_X(x) dx.\] We immediately observe a rather obvious problem. This integral is undefined at \(x = 0\) and diverges in a neighbourhood around zero. \citet{khuri2002} have stated the following corollary for the existence of a moment with negative first order:
\begin{corollary}
    If \(f_X(x)\) is a continuous pdf defined on \((-\infty, \infty)\), and if \[\lim_{x \to 0} \frac{f_X(x)}{|x|^\alpha} < \infty\] for \(\alpha > 0\), then \[\mathbb{E}[X^{-1}] \text{ exists}.\]
\end{corollary}

Most common distribution functions do not adhere to this corollary, however, the Gamma function does (see example \ref{p:negative}).

If such a moment of negative order exists, we should be able to obtain it using the MGF. In the 20-th century, \citet{cressie1981} have published the following remarkable theorem:

\begin{theorem}\label{t: negative}
    Assuming the negative \(n\)-th raw moment exists, the negative \(n\)-th raw moment can be computed as follows: 
    \[\mathbb{E}[X^{-n}] = \frac{1}{\Gamma(n)} \int_{0}^{\infty} t^{n- 1} M_X(-t) dt\] where \(n\) is a positive integer.
\end{theorem}
The proof of this Theorem can be found in Appendix \ref{pf:MGF}.

Since this is an extension on the regular functions of the MGF, this technique is of interest for this thesis. Thus, it will be shortly be discussed. We compute the first inverse moment of the Gamma distribution, by making use of the latter theorem for the MGF.

\begin{example}
    Let \[f_X(x) \sim \Gamma(\alpha, \lambda) = 
    \frac{x^{\alpha -1} e^{-\lambda x} \lambda^\alpha} {\Gamma(\alpha)}, M_X(t) = \left(\frac{\lambda}{\lambda - t}\right)^\alpha\]
    \[\mathbb{E}[X^{-1}] = \frac{1}{\Gamma(1)} \int_{0}^{\infty} t^{( 1 - 1)} \left(\frac{\lambda}{\lambda - (-t)}\right)^\alpha dt =  \int_{0}^{\infty} \left(\frac{\lambda}{\lambda + t}\right)^\alpha dt\]
    \[ = \lambda^\alpha \int_{0}^{\infty} (\lambda + t)^{-\alpha} dt, \text{ Let } u = \lambda + t, \frac{du}{dt} = 1, dt = du:\]
    \[ \lambda^\alpha \int_{0}^{\infty} u^{-\alpha} du
    =  \lambda^\alpha \frac{u^{ 1-\alpha}}{1 -\alpha}\Big|_{0}^{\infty} = \lambda^\alpha \frac{(\lambda + t)^{1 -\alpha}}{1 -\alpha}\Big|_{0}^{\infty}\]
    \[= \lambda^\alpha\left( 0 - \frac{\lambda^{ 1 - \alpha}}{1 -\alpha}\right) = \frac{-\lambda}{ 1 - \alpha} = \frac{\lambda}{\alpha - 1}.\] Which corresponds with our result from example \ref{p:negative}.
\end{example}

\subsubsection{Extending the Moment Generating Function to fractional order}
Now that we have introduced the definitions of the MGF and discussed a number of relevant properties, we will combine these with the techniques developed in section \ref{s:calculus}. Therefore, we can at last obtain moments of fractional order using the MGF.


To avoid confusion regarding what fractional derivative is being used in combination with the MGF, we will from now on, work with the following notation:
\begin{definition}\label{d: MGF}
    We define the MGF of order \(\alpha \in \mathbb{R}\) by \(\leftindex_{RL}{M}_X^{(\alpha)}, \leftindex_{CF}{M}_X^{(\alpha)}, \leftindex_{GL}{M}_X^{(\alpha)}\) for the MGF in combination with the Riemann-Liouville, Caputo-Fabrizio and Grünwald-Letnikov fractional derivative respectively.
\end{definition}
\begin{remark}
    The three properties mentioned in proposition \ref{p: moments} still hold for the MGF of fractional order. This is the case as the first property makes makes use of the derivative of order 0. Which has been defined to just be the original function itself as stated in the second property of proposition\ref{p: calculus}. The other two properties do not involve any derivatives of any order. Thus, they are generally applicable to the MGF, regardless of its order or kind of derivative.
\end{remark}

Before stating any results about the accuracy of these new MGF expressions, we first consider an explicit example to illustrate the interaction between fractional derivatives and the MGF. 

\begin{example}
    \begin{enumerate}[(i)]
        \item We let \(f_X(x_i) \sim Bernoulli(p)\), with \(\mathbb{P}(X = 1) = p\) and with associated MGF expression: \(M_X(t) = (1 - p) + p \cdot\exp(t)\), now we consider \[\leftindex_{CF}{M}_X^{(\frac{1}{2})} = \frac{1}{1 - \frac{1}{2}}  \int_{-h}^{t} \exp\left(\frac{-\frac{1}{2}}{1 - \frac{1}{2}}(t-s)\right) M_X'(s) ds.\] In this case, the domain of \(M_X(t) = (-\infty, \infty)\), so we let \(-h = -\infty\), and \(M_X'(s) = p\cdot \exp(s)\), thus we obtain:
        \[\leftindex_{CF}{M}_X^{(\frac{1}{2})} = 2  \int_{-\infty}^{t} \exp\left((s-t)\right) \cdot (p\cdot \exp(s)) ds.\]
        \[= 2p \cdot \exp(-t) \int_{-\infty}^{t}\exp(2s) ds\] 
        \[= p\cdot \exp(-t) \left(\exp(2s) \Big|_{-\infty}^{t}\right) = p\cdot \exp(t)\]
        Now, all that is left to do is set \(t = 0\) and we obtain that \(\leftindex_{CF}{M}_X^{(\frac{1}{2})} = p\).
        \item Computing \(\mathbb{E}[{X^{\frac{1}{2}}}]\) in the traditional fashion, we obtain: 
        \[\mathbb{E}[{X^{\frac{1}{2}}}] = \sum_x x^{\frac{1}{2}} \mathbb{P}(X = x) = 0^{\frac{1}{2}} \cdot \mathbb{P}(X = 0) + 1^{\frac{1}{2}} \cdot \mathbb{P}(X = 1)\]
        \[ = 0 \cdot(1 - p) + 1 \cdot p = p\]
    \end{enumerate}
    
\end{example}
It is amazing and maybe even somewhat surprising that both expressions obtain the same result. Indeed, this result is actually more of a coincidence. It is important to note that this agreement of results is coincidental and specific to the chosen distribution. Namely, all raw higher moments of a Bernoulli random variable are \(p\). If we had taken any other distribution in combination with a moment of fractional order, it becomes highly likely that the MGF returns a different value compared to the traditional method of computing moments. What is more, if we were to take \[\leftindex_{RL}{M}_X^{(\frac{1}{2})}  = \frac{d}{dt} \frac{1}{\sqrt{\pi}}  \int_{-h}^{t} (t - s)^{\frac{-1}{2}} f(s) ds\] we obtain an integral which may diverge based on the choice of \(-h\) . These observations lead to the following theorem.

\begin{theorem}\label{t: MGF_accurate}
     Consider the three MGF's as defined in definition \ref{d: MGF}. Assume \(\leftindex_{RL}{M}_X^{(\alpha)}\) and \(\leftindex_{GL}{M}_X^{(\alpha)}\) are well defined on some open interval \((-h, h)\), then the MGF expressions \(\leftindex_{RL}{M}_X^{(\alpha)}\) and \(\leftindex_{GL}{M}_X^{(\alpha)}\) accurately obtain raw moments of order \(\alpha \in \mathbb{R}\)
    
\end{theorem}
The proof can be found in Appendix \ref{pf:MGF}.
\newline
Unfortunately, this result does not hold for \(\leftindex_{CF}{M}_X^{(\alpha)}\), which leads to the following theorem.

\begin{theorem}\label{t: MGF_inaccurate}
    Consider the three MGF's as defined in definition \ref{d: MGF}. Assume \(\leftindex_{CF}{M}_X^{(\alpha)}\) is well defined on some open interval \((-h, h)\), then the MGF \(\leftindex_{CF}{M}_X^{(\alpha)}\) inaccurately approximates moments of order \(\alpha \in \mathbb{R}\) with approximation error given by
    \[
\begin{cases} 
    \displaystyle \int_{-\infty}^{\infty} x^\alpha  f_X(x) dx -  \displaystyle \int_{-\infty}^{\infty}  \frac{x^{n+1} }{(1 - \beta)x + \beta} f_X(x) dx & \text{if } X \text{ is continuous,} \\ 
    \displaystyle \sum_{i} \left(x_i^\alpha -  \frac{x_i^{n+1} }{(1 - \beta)x_i + \beta}\right) f_X(x_i) & \text{if } X \text{ is discrete.} 
\end{cases}
\] with \(\alpha \in \mathbb{R}, \beta = \alpha - n \text{ and } n = \lfloor \alpha \rfloor.\)
    
\end{theorem}
The proof can be found in Appendix \ref{pf:MGF}.

\begin{remark}
    In the case when \(\alpha \in \mathbb{N}\), we have that \(n = \alpha\), and thus \( \beta = 0\), therefore, \(\leftindex_{CF}{M}_X^{(\alpha)}\) is accurate for integer orders. This is an expected result, as the MGF for integer moments is accurate and from section \ref{s:calculus} we know that \(D_{CF}^{\alpha + \beta}f(x) = D_{CF}^\alpha(D_{CF}^\beta f(x))\), with \(\alpha \in \mathbb{N}, \beta \in [0, 1).\) In this case, let \(\beta = 0\). So we get \(D_{CF}^{\alpha + 0}f(x) = D_{CF}^{\alpha}f(x)\) which is just a regular derivative of integer order.
\end{remark}






\section{Conclusion}\label{s:con}
The conclusion summarizes the results of your research.
Also it gives possible directions of further research.

\section*{Acknowledgements}
Place the acknowledgments section, if needed, after the main text, 
but before any appendices and the references. The section heading is not numbered.

\newpage
\appendix

\section{Appendix: Relevant Functions and Identities}\label{s:appendices}

\begin{definition}\label{def:laplace}
    Let \(\displaystyle F(s) = \mathcal{L}(f) = \int_{0}^{\infty}e^{-st} f(t) dt\) be the Laplace transform of \(f\). Then \(f = \mathcal{L}^{-1}(f)\) is said to be the inverse Laplace transform of \(F\). The Laplace transform can be used to convert differential equations to much simpler algebraic equations \citep{thiagarajan2022}.
\end{definition}

\begin{definition}\label{def:mellin}
    The Mellin transform of some function \(f(t)\) is defined as \(\mathcal{M}\{f(t)\}(s) = \displaystyle \int_{0}^{\infty} t^{s-1} f(t) dt\). The Mellin transform is often used in computer science for the analysis of algorithms \citep{flajolet1996}.
\end{definition}

\begin{definition}\label{def:fourier}
    The Fourier transform of some function \(f(x)\) is defined as \(\displaystyle \hat{f}(\xi) = \int_{-\infty}^{\infty} f(x) e^{-2 \pi i \xi x}dx\), \(\xi \in \mathbb{R}\). The Fourier transform can be used to simplify differential equations and finds applications in signal and image processing \citep{fourier1822}. 
    
\end{definition}

\begin{definition}
    Let \(D\) be the differential operator, such that \(D f(x) = \frac{d}{dx} f(x)\). Then, as in \citet{samko1993}, the fractional derivative of order \(\alpha\) is defined as \[D^{\alpha} f(x) = \frac{d^{\alpha}}{dx^{\alpha}} f(x)\]
    
\end{definition}
In this definition, \(\alpha\) can be any real number. When taking regular derivatives, \(\alpha \in \mathbb{N}\). In most of our cases, we are interested in the instance where  \(\alpha \in \mathbb{R}_+\). It is also possible to study derivatives of negative order, which can be used to obtain moments of negative order of a random variable, provided that such an order exists. A derivative of negative order is simply an integral of positive order. This is defined as follows:
\begin{definition}
    Let \(I\) be the integral operator, such that \(I f(x) = \int f(x) dx\). Then by \citet{cauchy1823} the fractional integral of order \(\alpha\) is defined as \[(I^{\alpha} f) (x) = \frac{1}{(\alpha-1)!}\int (x-t)^{\alpha-1} f(t) dt\]
\end{definition} 

Combining the previous two definitions, we obtain the following, more general, definition, as in \citet{oldham1974}.
\begin{definition}\label{d: differintegral}
    The differintegral operator is defined as
    \begin{equation}
        R^\alpha f(x) = \begin{cases}
            I^{|\alpha|} f(x) & \text{if } \alpha < 0 \\
            D^\alpha f(x) & \text{if } \alpha > 0 \\
            f(x) & \text{if } \alpha = 0
        \end{cases}, \text{ with } \alpha \in \mathbb{R}.
        \end{equation}
        
\end{definition}

We define the (Euler-)Gamma function as follows:
\begin{definition}\label{d: eg}
    Let \(z \in \mathbb{C}\), for \(\Re(z) > 0\), we have the following: \(\displaystyle \Gamma(z) = \int_{0}^{\infty} t^{z-1} e^{-t} dt\)
\end{definition}

The Gamma function can be seen as an extension of the factorial function, for non-integers. This function is defined for complex numbers and all there subsets (so also real numbers), as long as the condition above holds. For positive integers values \(z\), we have the following identity: \(\Gamma(z) = (z - 1)!\)
Other important identities, not necessarily for \(z\) an integer, are: 
\begin{itemize}
    \item \(\Gamma(z + 1) = z \Gamma(z)\)
    \item \(\Gamma(2) = \Gamma(1) = 1\)
    \item \(\Gamma(\frac{1}{2}) = \sqrt{\pi}\)
\end{itemize}


\begin{definition}\label{def:binomial}
    For \(0 \leq k \leq n\), the Binomial Coefficient is defined as follows: \[\binom{n}{k} = \frac{n!}{k! (n - k)!} = \frac{\Gamma(n + 1)}{\Gamma(k + 1) \cdot \Gamma( n - k  + 1)}\] where \(n, k \in \mathbb{N}\) and \(\Gamma(.)\) as defined in \ref{d: eg}.
\end{definition}


\begin{definition}\label{def:vandermonde}
    Vandermonde's identity: for non-negative integers, \(k, l, m, n\), we have that \[\sum_{k = 0}^{l} \binom{m}{k} \cdot \binom{n}{l - k} = \binom{m + n}{l}\]
\end{definition}
A modification on the latter identity has been called the Chu-Vandermonde identity. This is the same identity, but it his been proven that the identities still hold for complex values \(m, n\) as long as \(l\) is a positive integer \citep{askey75}.

\begin{definition}\label{d: binomial}
    The binomial series is a generalization of the binomial formula, namely:
    \[(1 + x)^\alpha = \sum_{k = 0}^{\infty}\binom{\alpha}{k} \cdot x^k\]
\end{definition}

For the interchange-ability of derivatives and integrals and sums, we can apply the following two theorems:

\begin{theorem}\label{thm:leibnitz}
    Leibnitz's Rule ( as stated in \citet{flanders1973}):
    Let \(f(x, \theta), a(\theta), b(\theta)\) be differentiable with respect to \(\theta\), then we have that:
    \[ \frac{d}{d\theta} \int_{a(\theta)}^{b(\theta)} f(x, \theta) dx = f(b(\theta), \theta) \frac{d b(\theta)}{d\theta} - f(a(\theta), \theta) \frac{d a(\theta)}{d\theta} + \int_{a(\theta)}^{b(\theta)} \frac{\partial f(x, \theta) }{\partial \theta} dx.\] For the special case, where \(a(\theta), b(\theta)\) are constant we have that: 
    \[\frac{d}{d\theta} \int_{a}^{b} f(x, \theta) dx =  \int_{a}^{b} \frac{ \partial f(x, \theta)}{ \partial d\theta}.\]
\end{theorem}

For the interchange-ability of derivatives and summations, the following theorem has been given by \citet{casella2002}:
\begin{theorem}\label{thm:casella}
    Suppose that the series \(\sum_{x = 0}^{\infty} h(\theta, x)\) converges for all \(\theta\) in an interval \((a, b)\) of real numbers and 
    
    \begin{enumerate}[(i)]
        \item \(\frac{\partial h(\theta, x)}{\partial \theta}\) is continuous for all \(x\)
        \item \(\sum_{x = 0}^{\infty} \frac{\partial h(\theta, x)}{\partial \theta}\) converges uniformly on every closed bounded sub-interval of \((a, b)\)
    \end{enumerate}
    Then:
    \[
    \frac{d}{d \theta} \left( \sum_{x = 0}^{\infty} h(\theta, x) \right) = \sum_{x = 0}^{\infty} \frac{\partial h(\theta, x)}{\partial \theta}
    \]
    
   
\end{theorem}


\clearpage
\section{Appendix: Proofs of section \ref{s:methodology}}\label{s:app_B}

\subsection{Proofs section \ref{s:calculus}}\label{pf:calculus}

\begin{proof}
    Proof of \ref{p: calculus}
    \begin{enumerate}[(i)]
        \item We will prove for the Riemann-Liouville derivative, the proof for the Caputo-Fabrizio derivative is very similar and the Grünwald-Letnikov derivative is a direct consequence of the linearity of the sum.
        \[ D^{\alpha} (a f(x) + b g(x)) = \frac{d^n}{dx^n}\frac{1}{\Gamma(n - \alpha)} \int_{0}^{x} (x - t)^{n - \alpha - 1} \left(a f(t) + b g(t)\right) dt \]
     
        \[= \frac{d^n}{dx^n}\left(\frac{a}{\Gamma(n - \alpha)} \int_{0}^{x} (x - t)^{n - \alpha - 1}  f(t)dt + \frac{b}{\Gamma(n - \alpha)} \int_{0}^{x} (x - t)^{n - \alpha - 1} g(t)dt\right) \] Where we simply split the integral and put the constants in front.
        \[= \frac{d^n}{dx^n}\frac{a}{\Gamma(n - \alpha)} \int_{0}^{x} (x - t)^{n - \alpha - 1}  f(t)dt + \frac{d^n}{dx^n} \frac{b}{\Gamma(n - \alpha)} \int_{0}^{x} (x - t)^{n - \alpha - 1} g(t)dt \] As the regular derivative operator is linear.
        \[ = a D^{\alpha} f(x) + b D^{\alpha} g(x) \]
        \item Intuitively, this makes perfect sense, as the 0-th derivative is just no derivative, thus simply the function \(f(x)\). But for these derivatives, a little bit more effort is required to prove this rather obvious fact.
        \newline 
        For the Grünwald-Letnikov derivative we get: \[D^0 f(x) = \lim_{h \to 0} \frac{1}{h^0} \sum_{k=0}^\infty (-1)^k \binom{0}{k} f(x - k h)
        = \lim_{h \to 0} \frac{1}{1} \sum_{k=0}^\infty (-1)^k \frac{0!}{k!(0- k)!} f(x - k h).\] The factorial identity of the binomial coefficient only holds for \(0 \leq k \leq \alpha\). Since \(\alpha = 0\) and k is always a positive integer lesser or equal to \(\alpha, k = 0\). Thus, we get:
        \[ = \lim_{h \to 0} \sum_{k=0}^\infty (-1)^0 \frac{0!}{0!(0- 0)!} f(x - 0 h) = \lim_{h \to 0} f(x - 0 h) = f(x).\]
        \newline
        For the Caputo-Fabrizio derivative, we obtain the following:
        \[ D^{0} f(x) = \frac{1}{1 - 0}  \int_{0}^{x} \exp\left(\frac{0}{1 -0}(x-t)\right) f'(t) dt = \int_{0}^{x}f'(t) dt = f(x) .\]
        (By the first fundamental theorem of calculus).
        \newline
        Finally, for the Riemann-Liouville derivative, we can simply make use of \ref{r: integer} to note that in this context \(\alpha = 0\) is included in the natural integers. So \( \displaystyle D^\alpha = \frac{d^\alpha}{dx^\alpha} f(x) = \frac{d^0}{dx^0} f(x) = f(x)\). 

        \item The proof for the Riemann-Liouville derivative is given by \citet{koning15}. And the proof for the Caputo-Fabrizio derivative is given by \citet{losada15}. For the Grünwald-Letnikov derivative we get:
        \[ D^\alpha(D^\beta f(x)) = \lim_{h \to 0} \frac{1}{h^\alpha} \sum_{k=0}^\infty (-1)^k \binom{\alpha}{k} \left( \frac{1}{h^\beta} \sum_{l=0}^\infty (-1)^l \binom{\beta}{l} f(x - l h - kh)\right)\]
        \[= \lim_{h \to 0} \frac{1}{h^{\alpha + \beta}} \sum_{k=0}^\infty (-1)^k \binom{\alpha}{k} \sum_{l=0}^\infty (-1)^l \binom{\beta}{l} f(x - (k + l)h).\] We substitute \(m = k + l\) to deal with the double sums: 
        \[ \lim_{h \to 0} \frac{1}{h^{\alpha + \beta}} \sum_{m=0}^\infty f(x - mh)  \sum_{k=0}^m (-1)^k (-1)^{ m - k} \binom{\alpha}{k} \binom{\beta}{m - k}\] Now we make use of identify \ref{def:vandermonde} to obtain:
        \[ = \lim_{h \to 0} \frac{1}{h^{\alpha + \beta}} \sum_{m=0}^\infty (-1)^m \binom{\alpha + \beta}{m} f(x - mh) = D^{\alpha + \beta} f(x).\]
        It can be shown in a similar way that the latter expression is equal to \(D^\beta(D^\alpha f(x))\).

        
    \end{enumerate}
\end{proof}

\subsection{Proofs section \ref{s:MGF}} \label{pf:MGF}

\begin{proof}\label{p:negative}
    Example: let \(f_X(x) \sim \Gamma(\alpha, \lambda) = 
    \displaystyle \frac{x^{\alpha -1} e^{-\lambda x} \lambda^\alpha} {\Gamma(\alpha)}.\) This PDF is defined on \((0, \infty)\). So the function is not defined on \(\mathbb{R}\). This, however, is not a problem, as we can just evaluate the right limit. Since the Gamma function already uses \(\alpha\) as a parameter, we will evaluate \( \displaystyle \frac{f_X(x)}{|x|^\beta}, \text{ with }\beta > 0\):
    \[\lim_{x \to 0_+} \frac{f_X(x)}{|x|^\beta} = \lim_{x \to 0_+} \frac{x^{\alpha -1 - \beta} e^{-\lambda x} \lambda^\alpha} {\Gamma(\alpha)} = \lim_{x \to 0_+} \frac{x^{\alpha -(1 + \beta)} e^{-\lambda x} \lambda^\alpha} {\Gamma(\alpha)} .\] Thus for \(\alpha \geq \beta + 1, \lim_{x \to 0_+} \frac{f_X(x)}{|x|^\beta} < \infty\). Therefore, the first negative moment of the Gamma distribution should exist.

    We compute the first negative moment: 
    \[ \mathbb{E}[X^{-1}] = \int_{0}^{\infty} x^{-1} \frac{x^{\alpha -1} e^{-\lambda x} \lambda^\alpha} {\Gamma(\alpha)} dx = \frac{\lambda^\alpha}{\Gamma(\alpha)} \int_{0}^{\infty} x^{\alpha - 2} e^{-\lambda x} dx\]
    Using the substitution \(\displaystyle u = \lambda x, \frac{du}{dx} = \lambda, dx = \frac{du}{\lambda}\), we get:
    \[ = \frac{\lambda^\alpha}{\Gamma(\alpha)} \int_{0}^{\infty}\left(\frac{u}{\lambda}\right)^{\alpha -2} e^{-u} du = \frac{\lambda^\alpha}{\Gamma(\alpha) \lambda^{\alpha-1}} \int_{0}^{\infty}\left(\frac{u}{\lambda}\right)^{\alpha -2} e^{-u} du.\] This integral is equal to \(\Gamma(\alpha - 1)\) (See definition \ref{d: eg}). So we get: 
    \[\mathbb{E}[X^{-1}] = \frac{\lambda^\alpha \Gamma(\alpha - 1) }{\Gamma(\alpha) \lambda^{\alpha-1}} = \frac{\lambda \Gamma(\alpha - 1)}{(\alpha -1)\Gamma(\alpha - 1)} = \frac{\lambda}{(\alpha -1)}.\] Thus, for \(\alpha \neq 1, \mathbb{E}[X^{-1}] = \frac{\lambda}{(\alpha -1)}\). Fortunately, this is always the case, since we had just derived that the integral only converges when \(\alpha \geq \beta + 1, \text{ with } \beta > 0\). In other words, \(\alpha > 1\). So this holds.
\end{proof}

\begin{proof} Proof of Proposition \ref{t:mgf}

    \[M_X^{(n)}(t) = \frac{d^n}{dt^n} \int_{-\infty}^{\infty} e^{tx} f_X(x) dx = \int_{-\infty}^{\infty} \frac{d^n}{dt^n} e^{tx} f_X(x) dx\] (We can interchange differentiation and integration since all partial derivatives of \(e^{tx} f(x)\) are continuous and the absolute value of the integral converges, as we assume the \(n\)-th moment exists, see Theorem \ref{thm:leibnitz}).
    \[ = \int_{-\infty}^{\infty} x^n e^{tx} f_X(x) dx, \text{evaluating at } t = 0: = \int_{-\infty}^{\infty} x^n e^{0x} f_X(x) dx\]
    \[ = \int_{-\infty}^{\infty} x^n f_X(x) dx = \mathbb{E}[X^n]\]
    The proof for the case that \(f_X(x)\) is discrete is very similar. In that case, one would have to change the order of the derivative and summation, which has been justified in Theorem \ref{thm:casella}.
\end{proof}

\begin{proof}
    Proof of Proposition \ref{p: moments}
    \begin{enumerate}[(i)]
        \item This is trivial. For \(X\) a continuous random variable, we get: 
        \[ M_X^{(0)}(t) = \int_{-\infty}^{\infty} x^0 \cdot e^{t x } f_X(x) dx = \int_{-\infty}^{\infty} 1 \cdot e^{0 x } f_X(x) dx = \int_{-\infty}^{\infty} f_X(x) dx.\] Assuming that \(f(x)\) is a PDF, this integrates to 1 by definition. If this integral is not equal to 1, this implies that \(f(x)\) is not a PDF. The proof for the discrete case is the same but with a summation instead of an integral sign.
        \item \[M_{\mu + \sigma X}(t) = \mathbb{E}[e^{(\mu + \sigma X)t}] = \mathbb{E}[e^{\mu t} \cdot e^{\sigma X t}].\] Since this is the expectation of \(x\), every term that is not dependent on \(x\) can be taken out of the summation:
        \[ = e^{\mu t} \cdot \mathbb{E}[e^{\sigma X t}] = e^{\mu t} \cdot M_{ X}(\sigma t)\]
        \item 
        \[M_{X+Y}(t) = \mathbb{E}[e^{(X + Y)t}] = \int_{-\infty}^{\infty} \int_{-\infty}^{\infty} e^{(x + y)t} f_{X, Y}(x, y) dx dy\] 
        \[= \int_{-\infty}^{\infty} \int_{-\infty}^{\infty} e^{xt} \cdot e^{yt} f_{X, Y}(x, y) dx dy.\] 
        \(f_{X, Y}(x, y)\) is the joint pdf for \(X, Y\). But we know that the latter is equal to \(f_X(x) \cdot f_Y(y)\), if \(X, Y\) are independent. Thus we get:
        \[ = \left(\int_{-\infty}^{\infty}  e^{xt} f_X(x) dx\right) \left(\int_{-\infty}^{\infty}  e^{yt} f_Y(y) dy\right) = M_X(t) \cdot M_Y(t).\] 
    \end{enumerate}
\end{proof}

\begin{proof}
    Proof of Theorem \ref{t: negative}:
Suppose for the moment that \( X \) is a positive random variable. Since \( x \cdot f_X(x) \) is integrable for \( x > 0 \), we have:

\[
    \mathbb{E}(X) = \int_0^\infty x \, dF(x) = \int_0^\infty \int_{-\infty}^0 e^{tx} \, dt \, dF(x).
\]


\[
    =  \int_{-\infty}^0 \int_0^\infty e^{tx} \, dt \, dF(x) = \int_0^\infty M_X(-t) \, dt.
\]

The interchange of the order of integration is subject to \( \mathbb{E}(e^{-tX}) \) being integrable from \( t = 0 \) to \( t = \infty \).

Finally, by substituting \( X^{-1} \) for \( X \), we find:

\[
    \mathbb{E}(X^{-1}) = \int_0^\infty M_X^{-1}(-t) \, dt.
\]

There are two natural ways to generalize (1) to \( \mathbb{E}(X^{-1}) \); one way gives:

\[
    \mathbb{E}(X^{-n}) = \int_0^\infty \int_0^\infty \cdots \int_0^\infty M_X(-t_n) \, dt_n \cdots dt_2 dt_1, \tag{2}
\]

while the second way gives:

\[
    \mathbb{E}(X^{-n}) = \frac{1}{\Gamma(n)} \int_0^\infty t^{n-1} M_X(-t) \, dt.
\]
\citep{cressie1981}
\end{proof}

\begin{proof}
    Proof Theorem \ref{t: MGF_accurate}.
    We consider \(X\) to be a continuous variable. The three MGF expressions are accurate if \(\mathbb{E}[X^\alpha] - M_X^{(\alpha)}(0) = 0\), in other words, if 
    \[\int_{-\infty}^{\infty} (x^\alpha - D^\alpha e^{tx}) \cdot f_X(x) dx = 0 \iff x^\alpha = D^\alpha e^{tx}\]
    where \(D^\alpha\) the differential operator of order \(\alpha\). Note that we are taking derivatives w.r.t. \(t\), not \(x\). The proof will be shown for the Grünwald-left derivative.
    \begin{enumerate}[(i)]
        \item \[\leftindex_{GL}{M}_X^{(\alpha)} = D^\alpha_{GL}(e^{tx})  = \lim_{h \to 0} \frac{1}{h^\alpha} \sum_{k=0}^\infty (-1)^k \binom{\alpha}{k} \exp(x(t - kh))\]
        \[= \exp(xt) \lim_{h \to 0} \frac{1}{h^\alpha} \sum_{k=0}^\infty (-1)^k \binom{\alpha}{k} \exp(-xh)^k\]
        \[= \exp(xt) \lim_{h \to 0} \frac{1}{h^\alpha} (1 - \exp(-xh))^\alpha\]
        (where we used the binomial series identity defined in \ref{d: binomial}). Now we use the Taylor expansion of \(\exp(-xh) = 1 - xh + \frac{(xh)^2}{2!} - \frac{(xh)^3}{3!} + ... - ...\). Since \(h \to 0\), we get: \(\exp(-xh) = 1 - xh + \mathcal{O}(h^2)\) (the rest of the terms are negligible).
        Thus we get:
        \[ \exp(xt) \lim_{h \to 0} \frac{1}{h^\alpha} \left(1 -(1 - xh + \mathcal{O}(n))\right)^\alpha = \exp(xt) \lim_{h \to 0} \frac{1}{h^\alpha} \left(h(x + \mathcal{O}(h))\right)^\alpha \]
        \[ = \exp(xt) \lim_{h \to 0} \frac{1}{h^\alpha} \left(h^\alpha((x + \mathcal{O}(h)))^\alpha\right) = \exp(xt) \lim_{h \to 0} (x + \mathcal{O}(h))^\alpha\]
        \[ = \exp(xt) \cdot x^\alpha.\]
        Finally, we take a value of \(t\) around 0 and obtain \(\leftindex_{GL}{M}_X^{(\alpha)}(0) = x^\alpha\) 
        \item We consider the Riemann-Liouville derivative: 
        \[D_{RL}^\alpha(e^{tx}) = \frac{d^n}{dt^n} \frac{1}{\Gamma(n -\alpha)}  \int_{-\infty}^{t} (t-s)^{n - \alpha-1} \cdot e^{sx} ds\]
        \citet{koning15} has shown that this derivative is equal to \(x^\alpha e^{tx}\). Thus, if we take a value of \(t\) around 0, we get: \(\leftindex_{RL}{M}_X^{(\alpha)}(0) = x^\alpha\)
    \end{enumerate}

    The proof also holds for the case when \(X\) is a discrete random variable.
\end{proof}

\begin{proof}
    Proof of Theorem \ref{t: MGF_inaccurate}.
    We compute \[\leftindex_{CF}{M}_X^{(\alpha)} = D^\alpha_{CF}(e^{tx}) = \frac{d^n}{dt^n}\frac{1}{1 - \beta}  \int_{-\infty}^{t} \exp\left(\frac{-\beta}{1 - \beta}(t - s)\right) x \exp(xs) ds\]
    \[= \frac{d^n}{dt^n}\frac{x}{1 - \beta} \exp\left(\frac{-\beta t}{1 - \beta}\right) \int_{-\infty}^{t} \exp\left(s\left(\frac{\beta }{1 - \beta} + x \right)\right) ds  \text{ where } \beta = \alpha - n \text{ and } n = \lfloor \alpha \rfloor.\]
    Now let \( \displaystyle u = \frac{\beta}{1 - \beta} + x\), so we get:
    \[\frac{d^n}{dt^n}\frac{x}{1 - \beta} \exp\left(\frac{-\beta t}{1 - \beta}\right) \int_{-\infty}^{t} \exp(us) ds = \frac{d^n}{dt^n}\frac{x}{1 - \beta} \exp\left(\frac{-\beta t}{1 - \beta}\right) \frac{1}{u} exp(us)\Big|_{-\infty}^{t}\]
    \[= \frac{d^n}{dt^n} \frac{x}{1 - \beta} \exp\left(\frac{-\beta t}{1 - \beta}\right)  \frac{1}{u}  \exp(u t) = \frac{d^n}{dt^n}\frac{x \exp(xt)}{(1 - \beta)x + \beta}\]
    Now we apply the derivative of the above expression, with respect to \(t\), \(n\) times:
    \[\frac{d^n}{d t^n} \frac{x \exp(xt)}{(1 - \beta)x + \beta} = \frac{x^{n+1} \exp(xt)}{(1 - \beta)x + \beta}\]
    Now we take a value of \(t\) around 0 and obtain:
    \[\leftindex_{CF}{M}_X^{(\alpha)} = D^\alpha_{CF}(e^{tx}) = \frac{x^{n+1} }{(1 - \beta)x + \beta}.\]
    The latter expression is not equal to \(x^\alpha\), thus the error of \(\leftindex_{CF}{M}_X^{(\alpha)}\) is:
    \[
    \begin{cases} 
    \displaystyle \int_{-\infty}^{\infty} x^\alpha  f_X(x) dx -  \displaystyle \int_{-\infty}^{\infty}  \frac{x^{n+1} }{(1 - \beta)x + \beta} f_X(x) dx & \text{if } X \text{ is continuous,} \\ 
    \displaystyle \sum_{i} \left(x_i^\alpha -  \frac{x_i^{n+1} }{(1 - \beta)x_i + \beta}\right) f_X(x_i) & \text{if } X \text{ is discrete.} 
\end{cases}\]

\end{proof}

\clearpage

\section{Appendix: Algorithms}\label{s:app_alg}
\begin{center}
\begin{minipage}{0.9 \textwidth}
\begin{algorithm}[H]
\caption{Compute Fractional Moment via CMGF Method}
\label{alg:cmgf_dist}
\begin{algorithmic}[1]
\REQUIRE Order $r$, location parameter $\xi$, scale parameter $s$, distribution name \texttt{dist\_name}, distribution parameters \texttt{parameters}, tolerance \texttt{tol}
\ENSURE Approximate value of the fractional moment
\STATE Define integrand function:
\STATE \hspace{1em} For $t \in [0, \infty)$:
\STATE \hspace{2em} $z \leftarrow s + it$
\STATE \hspace{2em} $Mz \leftarrow \texttt{MGF\_DIST}(z, \texttt{parameters}, \texttt{dist\_name})$
\STATE \hspace{2em} $M_{-z} \leftarrow \texttt{MGF\_DIST}(-z, \texttt{parameters}, \texttt{dist\_name})$
\STATE \hspace{2em} $\text{numerator} \leftarrow e^{-\xi z} \cdot Mz + e^{\xi z} \cdot M_{-z}$
\STATE \hspace{2em} $\text{integrand}(t) \leftarrow \Re\left(\dfrac{\text{numerator}}{z^{r+1}}\right)$
\STATE Compute integral: $\text{integral} \leftarrow \int_0^{100} \text{integrand}(t)\,dt$ using \texttt{quadgk} with relative tolerance \texttt{tol}
\STATE \textbf{return} $\dfrac{\Gamma(r+1)}{\pi} \cdot \text{integral}$
\end{algorithmic}
\end{algorithm}
\end{minipage}
\end{center}

\begin{center}
\begin{minipage}{0.9 \textwidth}
\begin{algorithm}[H]
\caption{Analyze Fractional Moments}
\label{alg:analyse_fractional_moments}
\begin{algorithmic}[1]
\REQUIRE Parameters: $\omega$, $\alpha$, $\beta$, list of horizons: $H\_values$, list of orders: $orders$, number of iterations: $n\_sim$
\ENSURE DataFrame of conditional moments for each $H$ and order
\STATE Initialize empty results DataFrame
\FOR{each $H$ in $H\_values$}
    \STATE $R_{H}^{2} \leftarrow$ SimulateConditionalVariance($\omega, \alpha, \beta, H, n\_sim$)
    \STATE $pdf\_R \leftarrow$ KDE distribution from $R_H^2$
    \FOR{each $order$ in $orders$}
        \STATE $empirical \leftarrow \text{mean}(|R_H^2|^{order})$
        \STATE $standard \leftarrow$ FractionalMoment($pdf\_R$, $order$, use\_abs=true)
        \STATE $CF \leftarrow$ FractionalMomentCF($pdf\_R$, $order$, use\_abs=true)
        \STATE Append $(H, order, empirical, standard, CF)$ to results
    \ENDFOR
\ENDFOR
\RETURN results
\end{algorithmic}
\end{algorithm}
\end{minipage}
\end{center}

\begin{center}
\begin{minipage}{0.9 \textwidth}
\begin{algorithm}[H]
\caption{Simulate Conditional Variance}
\label{alg:simulate_conditional_variance}
\begin{algorithmic}[1]
\REQUIRE Parameters: $\omega$, $\alpha$, $\beta$, list of horizons: $H\_values$, number of iterations: $n\_sim$
\ENSURE Array $R_H^2$ of accumulated conditional variances
\STATE Initialize $R_H^2$ as a zero array of length $n\_sim$
\STATE $\sigma_0^2 \leftarrow \omega / (1 - \alpha - \beta)$
\FOR{$i = 1$ to $n\_sim$}
    \STATE $\sigma^2 \leftarrow \sigma_0^2$
    \STATE $variance\_R_H \leftarrow 0$
    \FOR{$h = 1$ to $H$}
        \STATE $\sigma \leftarrow \sqrt{\sigma^2}$
        \STATE $r \sim \mathcal{N}(0, \sigma)$
        \STATE $variance\_R_H \leftarrow variance\_R_H + r^2$
        \STATE $\sigma^2 \leftarrow \omega + \alpha r^2 + \beta \sigma^2$
    \ENDFOR
    \STATE $R_H^2[i] \leftarrow variance\_R_H$
\ENDFOR
\RETURN $R_H^2$
\end{algorithmic}
\end{algorithm}
\end{minipage}
\end{center}


\begin{center}
\begin{minipage}{0.9 \textwidth}
\begin{algorithm}[H]
\caption{Compute $z$ from returns using GARCH and fractional moments}
\label{alg:compute_z}
\begin{algorithmic}[1]
\REQUIRE Parameters: $H \gets 2$, $n_{\text{sim}} \gets 100{,}000$, $\text{window} \gets 520$, $\text{target\_length} \gets$ \texttt{nothing}, returns, fractional order $\gamma$
\ENSURE Computed array \texttt{z}
\STATE 
\STATE $n \gets \text{length}(\texttt{returns})$
\STATE $\text{max\_index} \gets n - H$
\IF{$\text{target\_length} = \texttt{nothing}$}
    \STATE $\text{range} \gets \text{window} : \text{max\_index}$
\ELSE
    \STATE $\text{range} \gets (\text{max\_index} - \text{target\_length} + 1) : \text{max\_index}$
\ENDIF
\STATE Initialize empty array $z$
\FOR{each $t$ in \texttt{range}}
    \STATE $r_{\text{window}} \gets \texttt{returns}[(t - \text{window} + 1) : t]$
    \STATE Fit GARCH(1,1) model to $r_{\text{window}}$
    \STATE Extract parameters $\omega, \alpha, \beta$ from model
    \STATE $R_{H}^{2} \gets \text{simulate\_conditional\_variance}(\omega, \alpha, \beta, H, n_{\text{sim}})$
    \STATE $\text{distribution}_R \gets \text{kde}(R_{H}^{2})$
    \STATE $\text{pdf}_R \gets \text{KDEDistribution}(\text{distribution}_R)$
    \STATE $\text{moment} \gets \text{fractional\_moment}(\text{pdf}_R, \gamma; \text{use\_abs} = \text{true})$
    \STATE Append \texttt{moment} to \texttt{z}
\ENDFOR
\RETURN $z$
\end{algorithmic}
\end{algorithm}
\end{minipage}
\end{center}

\clearpage

\section{Appendix: Table of Common Distributions}\label{s:app_common_distributions}
\begin{table}[ht]
\renewcommand{\arraystretch}{1.5}
\begin{adjustwidth}{-.5in}{-.5in}
\begin{center}
    
\begin{tabular}{|c|c|c|c|}
\hline
\textbf{Distribution} & \textbf{PDF / PMF} & \textbf{MGF \(M_X(t)\)} & \textbf{Restrictions} \\
\hline

\( \text{Bernoulli}(p) \) & \(\displaystyle  f(x) = p^x(1-p)^{1-x} \) & \( M_X(t) = (1 - p) + pe^t \) & \( 0 \leq p \leq 1,\ x \in \{0,1\} \) \\
\hline

\( \text{Binomial}(n, p) \) & \( \displaystyle f(x) = \binom{n}{x} p^x (1-p)^{n-x} \) & \( M_X(t) = (1 - p + pe^t)^n \) & \( n \in \mathbb{N},\ 0 \leq p \leq 1 \) \\
\hline

\( \text{Poisson}(\lambda) \) & \( \displaystyle f(x) = \frac{\lambda^x e^{-\lambda}}{x!} \) & \( M_X(t) = \exp(\lambda(e^t - 1)) \) & \( \lambda > 0,\ x \geq 0 \) \\
\hline

\( \text{Exponential}(\beta) \) & \( \displaystyle f(x) = \frac{1}{\beta} e^{-x/\beta} \) & \( M_X(t) = \frac{1}{1 - \beta t},\ t < 1/\beta \) & \( \beta > 0,\ x > 0 \) \\
\hline

\( \text{Gamma}(\alpha, \beta) \) & \( \displaystyle f(x) = \frac{1}{\Gamma(\alpha)\beta^\alpha} x^{\alpha-1} e^{-x/\beta} \) & \( M_X(t) = (1 - \beta t)^{-\alpha},\ t < 1/\beta \) & \( \alpha, \beta > 0,\ x > 0 \) \\
\hline

\( \mathcal{N}(\mu, \sigma^2) \) & \( \displaystyle f(x) = \frac{1}{\sqrt{2\pi \sigma^2}} e^{-\frac{(x - \mu)^2}{2\sigma^2}} \) & \( M_X(t) = \exp\left( \mu t + \frac{\sigma^2 t^2}{2} \right) \) & \( \sigma > 0 \) \\
\hline

\end{tabular}
\end{center}
\end{adjustwidth}
\caption{Common distributions and their moment generating functions}
\label{t:MGFs_common}
\end{table}

\section{Appendix}\label{s:app1}
Place any appendices after the acknowledgments, starting on a new page.
The appendices are numbered
\textbf{A}, \textbf{B}, \textbf{C}, and so forth.
The appendices contain material that you would like to share with the
reader but that would hinder the flow of reading. For instance long
proofs of theorems, code of algorithms, data, etc.


%\printbibliography
% or in case of classic bibtex:
\bibliographystyle{chicago}
\bibliography{bib/bibfile}

\end{document}
